\newpage



\begin{define}{Permutations without Repetition}
Arranging $r$ distinct objects from $n$ distinct objects, where each object 
can be used at most once. Order matters, repetition not allowed.

$$P(n,r) = \frac{n!}{(n-r)!}$$
where $0 \leq r \leq n$.
\end{define}

\textbf{Example 1: Committee Leadership} (5 points)\\
A committee of 8 people needs to elect a president, vice president, and secretary. How many different ways can these three positions be filled?
\begin{subanswer}
    $P(8,3) = \frac{8!}{(8-3)!} = \frac{8!}{5!} = 8 \times 7 \times 6 = 336$ ways. We choose 3 distinct people from 8 and arrange them in order.
\end{subanswer}

\textbf{Example 2: Password Creation} (5 points)\\
How many 4-digit passwords can be created using the digits 1-9 if no digit can be repeated?
\begin{subanswer}
    $P(9,4) = \frac{9!}{(9-4)!} = \frac{9!}{5!} = 9 \times 8 \times 7 \times 6 = 3024$ passwords.
\end{subanswer}

\textbf{Example 3: Race Finishers} (5 points)\\
In a race with 12 participants, how many different ways can the top 5 finishers be ordered?
\begin{subanswer}
    $P(12,5) = \frac{12!}{(12-5)!} = \frac{12!}{7!} = 12 \times 11 \times 10 \times 9 \times 8 = 95,040$ ways.
\end{subanswer}

\begin{define}{Permutations with Repetition}
Arranging $n$ objects with $n_1$ identical of type 1, $n_2$ identical of type 2, 
..., $n_k$ identical of type $k$. Order matters, identical objects indistinguishable.

$$\frac{n!}{n_1! \cdot n_2! \cdot \cdots \cdot n_k!}$$
where $n = n_1 + n_2 + \cdots + n_k$.
\end{define}

\textbf{Example 4: Letter Arrangements} (5 points)\\
How many different arrangements can be made using all the letters in the word "MISSISSIPPI"?
\begin{subanswer}
    The word has 11 letters total: 1 M, 4 I's, 4 S's, and 2 P's. The number of arrangements is $\frac{11!}{1! \cdot 4! \cdot 4! \cdot 2!} = \frac{39,916,800}{1 \cdot 24 \cdot 24 \cdot 2} = 34,650$ arrangements.
\end{subanswer}

\textbf{Example 5: Coin Arrangements} (5 points)\\
How many different ways can 3 pennies, 2 nickels, and 1 dime be arranged in a row?
\begin{subanswer}
    Total of 6 coins: 3 pennies, 2 nickels, 1 dime. Number of arrangements = $\frac{6!}{3! \cdot 2! \cdot 1!} = \frac{720}{6 \cdot 2 \cdot 1} = 60$ ways.
\end{subanswer}

\textbf{Example 6: Binary Strings} (5 points)\\
How many different 8-bit binary strings can be formed with exactly 5 ones and 3 zeros?
\begin{subanswer}
    $\frac{8!}{5! \cdot 3!} = \frac{40320}{120 \cdot 6} = 56$ different binary strings.
\end{subanswer}


\newpage 



\begin{define}{Combinations without Repetition}
    Selecting $r$ objects from $n$ distinct objects, where order doesn't matter 
    and each object can be chosen at most once.
    
    $$\binom{n}{r} = \frac{n!}{r!(n-r)!}$$
    where $0 \leq r \leq n$.
  \end{define}

\textbf{Example 7: Committee Selection} (5 points)\\
A club has 12 members. How many different 5-member committees can be formed?
\begin{subanswer}
    $\binom{12}{5} = \frac{12!}{5!(12-5)!} = \frac{12!}{5! \cdot 7!} = \frac{479,001,600}{120 \cdot 5040} = 792$ committees.
\end{subanswer}

\textbf{Example 8: Lottery Numbers} (5 points)\\
In a lottery, players choose 6 numbers from 1 to 49. How many different combinations are possible?
\begin{subanswer}
    $\binom{49}{6} = \frac{49!}{6!(49-6)!} = \frac{49!}{6! \cdot 43!} = 13,983,816$ different combinations.
\end{subanswer}

\textbf{Example 9: Pizza Toppings} (5 points)\\
A pizza place offers 8 different toppings. How many different 3-topping pizzas can be ordered?
\begin{subanswer}
    $\binom{8}{3} = \frac{8!}{3!(8-3)!} = \frac{8!}{3! \cdot 5!} = \frac{40320}{6 \cdot 120} = 56$ different pizzas.
\end{subanswer}

\begin{define}{Combinations with Repetition}
    Distributing $n$ identical objects into $m$ distinct categories. 
    Order doesn't matter, objects can be reused. Equivalent to counting 
    non-negative integer solutions to $x_1+x_2+\cdots+x_m=n$.
    
    $$\binom{n+m-1}{n}$$
    where $x_i \geq 0$ for all $i \in\{1,2, \ldots, m\}$.
  \end{define}

\textbf{Example 10: Ice Cream Cones} (5 points)\\
An ice cream shop offers 5 flavors. How many different 3-scoop cones can be ordered if multiple scoops of the same flavor are allowed?
\begin{subanswer}
    $\binom{3+5-1}{3} = \binom{7}{3} = \frac{7!}{3!(7-3)!} = \frac{7!}{3! \cdot 4!} = \frac{5040}{6 \cdot 24} = 35$ different cones.
\end{subanswer}

\textbf{Example 11: Candy Distribution} (5 points)\\
How many ways can 10 identical candies be distributed among 4 children if each child can receive any number of candies (including 0)?
\begin{subanswer}
    $\binom{10+4-1}{10} = \binom{13}{10} = \binom{13}{3} = \frac{13!}{3!(13-3)!} = \frac{13!}{3! \cdot 10!} = \frac{6,227,020,800}{6 \cdot 3,628,800} = 286$ ways.
\end{subanswer}

\textbf{Example 12: Variable Solutions} (5 points)\\
How many non-negative integer solutions does the equation $x + y + z = 8$ have?
\begin{subanswer}
    $\binom{8+3-1}{8} = \binom{10}{8} = \binom{10}{2} = \frac{10!}{2!(10-2)!} = \frac{10!}{2! \cdot 8!} = \frac{3,628,800}{2 \cdot 40320} = 45$ solutions.
\end{subanswer}


\newpage

\begin{define}{Combinations with Repetition}
    The equation $x_1+x_2+\cdots+x_m=n$ has $\binom{n+m-1}{n}$ integer solutions such that $x_i \geq 0$ for any $i \in\{1,2, \ldots, m\}$.
  \end{define}
  
  
  
  
  \textbf{Strings with Multiple Restrictions: Permutations, Combinations, and Rule of Product} (10 points)\\
    Let $X$ be the set of the 26 lowercase English letters and 10 decimal digits. How many $X$-strings of length 15 satisfy all of the following properties (at the same time)?
    \begin{itemize}
      \item The first and last symbols of the string are distinct digits (which may appear elsewhere in the string).
      \item Precisely four of the symbols in the string are the letter 't'.
      \item Precisely three characters in the string are elements of the set $V=\{a,e,i,o,u\}$ and these characters are all distinct.
    \end{itemize}
    \begin{subanswer}
      $P(10,2) \cdot \binom{13}{4} \cdot \left(\binom{9}{3} \cdot P(5,3)\right) \cdot 30^6$. Place 2 distinct decimal digits at the first and last of the string, there are $P(10,2)$ such strings. Choose where to place the letter 't'. Since the first and last positions are not available, 4 positions will be chosen from 13 available positions: $\binom{13}{4}$. Choose 3 positions for vowels ($a,e,i,o,u$) from the rest of 9 positions, then place 3 distinct vowels in the positions: $\binom{9}{3} \cdot P(5,3)$. Lastly, we have 6 positions left to fill, and those 6 positions cannot contain vowels or 't' chosen previously: $30^6$.
    \end{subanswer}
  
  
  